\section*{How the Proof Works}

Take one vortex tube and stretch it along its axis. Incompressibility does not
allow that extension to create volume, so the tube contracts across its width.
The rotating core narrows. Velocity changes across a shorter distance, the
gradient steepens, and the vorticity becomes more concentrated even though the
total kinetic energy can remain finite. This is the danger: the energy does not
have to become infinite for the fluid to develop finer and more violent
structure.

Navier--Stokes has a scaling that follows this concentration exactly. If
\((u,p)\) is a solution, then
\[
u_\lambda(x,t)=\lambda u(\lambda x,\lambda^2t),
\qquad
p_\lambda(x,t)=\lambda^2p(\lambda x,\lambda^2t)
\]
preserves the balance between transport, pressure, and viscosity. Every term in
the equation acquires the factor \(\lambda^3\). Meanwhile,
\[
\|u_\lambda(\cdot,t)\|_{\dot H^s}
=\lambda^{s-1/2}
  \|u(\cdot,\lambda^2t)\|_{\dot H^s}.
\]
Concentration can make a norm below \(s=\tfrac12\) smaller and a norm above it
larger. The height \(s=\tfrac12\) does not move at all. It is the critical
height at which zooming into the narrowing core leaves the size of the
measurement unchanged.

That scaling gives the short Zeno calculation. Suppose one turn of the
cascade contracts the spatial scale by a fixed factor \(0<\rho<1\). A rescaled
copy then has
\[
r_j=\rho^j r_0,
\qquad
U_j=\rho^{-j}U_0,
\qquad
\tau_j=\rho^{2j}\tau_0.
\]
The velocity scale rises as the core narrows, while the parabolic response time
falls with the square of its width. Thus
\[
\sum_{j=0}^{\infty}\tau_j
=\tau_0\sum_{j=0}^{\infty}\rho^{2j}
=\frac{\tau_0}{1-\rho^2}<\infty.
\]
Scaling permits infinitely many finer turns to fit inside a finite interval.
This calculation does not produce a singularity. It shows why finite energy
and the parabolic scaling do not rule one out.

In order for a singularity to form, velocity has to blow up. In order for that
to happen, its acceleration would also have to blow up, which means its jerk
has to blow up, and so on and so on:
\[
u,\qquad
\partial_tu,\qquad
\partial_t^2u,\qquad
\partial_t^3u,\qquad\ldots
\]
What you'd expect to see in this pattern, if a singularity were to form, as
you go up the time derivatives, you'd expect to see less spatial regularity,
not more. Each higher temporal response should carry a rougher spatial field
as the cascade approaches its terminal time.

Now differentiate Navier--Stokes and watch what actually happens. Before any
shorthand is introduced, the first three equations are
\[
\begin{aligned}
\partial_tu
={}&\nu\Delta u
 -(u\!\cdot\nabla)u
 -\nabla p,\\[3pt]
\partial_t^2u
={}&\nu\Delta\partial_tu
 -(\partial_tu\!\cdot\nabla)u
 -(u\!\cdot\nabla)\partial_tu
 -\nabla\partial_tp,\\[3pt]
\partial_t^3u
={}&\nu\Delta\partial_t^2u
 -(\partial_t^2u\!\cdot\nabla)u
 -2(\partial_tu\!\cdot\nabla)\partial_tu
 -(u\!\cdot\nabla)\partial_t^2u
 -\nabla\partial_t^2p.
\end{aligned}
\]
The nonlinear term opens once, then twice, and its multiplicities begin to
appear. The complete pattern is the binomial form
\[
\partial_t^{\,n+1}u
=\nu\Delta\partial_t^nu
-\sum_{a=0}^{n}\binom na
  \bigl(\partial_t^au\!\cdot\nabla\bigr)
  \partial_t^{\,n-a}u
-\nabla\partial_t^np.
\]
Only now set
\[
V_n:=\partial_t^nu,
\qquad
Q_n:=\partial_t^np.
\]
The recurrence becomes
\[
V_{n+1}
+\sum_{a=0}^{n}\binom na
  (V_a\!\cdot\nabla)V_{n-a}
+\nabla Q_n
=\nu\Delta V_n,
\]
and pressure remains attached to every row through
\[
-\Delta Q_n
=\sum_{a=0}^{n}\binom na
  \partial_i\partial_j
  \bigl((V_a)_i(V_{n-a})_j\bigr).
\]
The term \(\nu\Delta V_n\) places two spatial derivatives of one response
inside the equation for the next temporal response. Transport and pressure
rise through this expansion, so none of those terms can be dropped when
the row is controlled.

At the critical height, the temporal climb exposes the spatial climb directly.
Let \(\Lambda=(-\Delta)^{1/2}\), and write
\[
E_s(t):=\frac12\|\Lambda^su(t)\|_2^2,
\qquad
H(t):=E_{1/2}(t).
\]
Pairing the equation at height \(s\) gives
\[
E_s'=\mathcal P_s-2\nu E_{s+1},
\]
where \(\mathcal P_s\) is the complete nonlinear transfer at that height. At
\(s=\tfrac12\), repeated differentiation makes the first steps visible:
\[
\begin{aligned}
H'
&=\mathcal P_{1/2}-2\nu E_{3/2},\\
H''
&=\partial_t\mathcal P_{1/2}
  -2\nu\mathcal P_{3/2}
  +(2\nu)^2E_{5/2},\\
H'''
&=\partial_t^2\mathcal P_{1/2}
  -2\nu\partial_t\mathcal P_{3/2}
  +(2\nu)^2\mathcal P_{5/2}
  -(2\nu)^3E_{7/2}.
\end{aligned}
\]
At every finite order this relation reads
\[
H^{(n)}
=(-2\nu)^nE_{n+1/2}
+\sum_{j=0}^{n-1}(-2\nu)^j
  \partial_t^{\,n-1-j}\mathcal P_{j+1/2},
\]
so the highest spatial energy can be isolated:
\[
E_{n+1/2}
=(-2\nu)^{-n}
\left(
H^{(n)}
-\sum_{j=0}^{n-1}(-2\nu)^j
 \partial_t^{\,n-1-j}\mathcal P_{j+1/2}
\right).
\]
The expected direction has reversed. Higher temporal order exposes higher
spatial Sobolev order. As you do that, you get higher order Sobolev embeddings.
For each finite \(k\), choosing \(n>k+1\) gives
\[
H^{n+1/2}(\mathbb R^3)\hookrightarrow C_b^k(\mathbb R^3),
\]
and carrying every finite height to a common endpoint gives the full payoff
\[
\bigcap_{m<\infty}H^m(\mathbb R^3)
=H^\infty(\mathbb R^3)
\hookrightarrow C^\infty(\mathbb R^3).
\]
This is one realization: the temporal sequence that should reveal increasing
spatial roughness has opened the route to every finite spatial height.

The recurrence is an identity, not an estimate. It identifies the higher
energy inside each temporal response, but it does not keep the nonlinear
transfers finite. Actual control enters one finite piece at a time. Fix a
finite temporal depth \(N\) and a finite spatial height \(M\). The
datum-generated whole-terminal rectangle theorem gives
\[
\sup_{0<T<T_*}\mathcal R_{N,M}(u;T)<\infty,
\]
with a constant allowed to depend on \(N\) and \(M\). No common bound over all
heights is being assumed. In particular, for every \(0\le n\le N\), the
rectangle contains the adjacent pair
\[
V_n\in L^2(0,T_*;H^{M+1}),
\qquad
V_{n+1}\in L^2(0,T_*;H^{M-1}).
\]
The adjacent rows give the temporal trace
\[
V_n\in C([0,T_*];H^M),
\]
and the pressure equation supplies the corresponding \(Q_n\)-coordinates.
This is a finite pressure-complete rectangle: finitely many temporal responses,
finitely many spatial heights, their terminal traces, and the pressure required
by the equation.

Now enlarge the rectangle. A larger choice of \((N,M)\) does not construct new
derivatives on the coordinates already present. Every rectangle is a
restriction of the datum-generated maximal history, so two rectangles agree
where they overlap. Their compatible projective intersection is
\[
\mathcal I[u_0]
=\bigl(\mathcal R_{N,M}[u_0;T_*]\bigr)_{N,M\in\mathbb N_0},
\qquad
u\in\bigcap_{r,m<\infty}
H^r(0,T_*;H^m(\mathbb R^3)).
\]
Each finite coordinate was controlled before it entered this intersection, and
compatibility makes all of those coordinates parts of one velocity--pressure
history.

Take its zeroth temporal row. For every finite \(m\), it has one trace in
\(H^m\), and agreement on the overlaps makes those traces one terminal
velocity rather than a separate endpoint for each height:
\[
u\in\bigcap_{m<\infty}C([0,T_*];H^m_\sigma),
\qquad
u_*:=u(T_*)
\in H^\infty_\sigma(\mathbb R^3)
\hookrightarrow C^\infty(\mathbb R^3).
\]
The recurrence and the pressure equation can now be evaluated at \(T_*\),
recovering every finite temporal and pressure coordinate of the terminal jet.
Local well-posedness starts a smooth solution from \(u_*\). On the interval
where that restart meets the original history, uniqueness identifies the two,
so the history continues past \(T_*\). If \(T_*\) had been a finite maximal
time, it was not maximal after all. Hence \(T_*=\infty\).

The familiar finite estimates now come from selecting the coordinate they need.
For finite temporal order \(r\), spatial derivative order \(k\), and
\(2\le p\le\infty\), choose a finite \(m\) with
\[
m>k+\frac32-\frac3p.
\]
On every finite time interval, the compatible intersection gives
\[
\mathcal I[u_0]
\longrightarrow H_t^1H_x^m
\longrightarrow C_tH_x^m
\longrightarrow L_t^qW_x^{k,p},
\qquad 1\le q\le\infty,
\]
for the chosen response \(V_r=\partial_t^ru\). Different finite choices give,
in immediate succession,
\[
\begin{gathered}
\mathcal I[u_0]\longrightarrow C_tH_x^1
\longrightarrow L_t^\infty L_x^6,\\[2mm]
\mathcal I[u_0]\longrightarrow C_tH_x^2
\longrightarrow L_t^\infty L_x^\infty,\\[2mm]
\mathcal I[u_0]\longrightarrow C_tH_x^3
\longrightarrow L_t^\infty W_x^{1,\infty}.
\end{gathered}
\]
No one of these estimates carries the argument. Any finite downstream demand
selects a sufficiently high finite coordinate from the compatible
intersection.
